% Canonical LaTeX source; imported and revised from the legacy Markdown.
\chapter{Numeri complessi}
\label{chap:5}

\begin{obiettivocapitolo}{scegliere la rappresentazione complessa più efficiente}
\prioritamedia. La forma algebrica è adatta a somme e parti reale/immaginaria; la forma esponenziale è adatta a prodotti, quozienti, potenze e radici.
\end{obiettivocapitolo}
\begin{proceduraesame}{potenze e radici complesse}
\begin{enumerate}
\item Scrivi $z=re^{i\theta}$ con $r=|z|$ e un argomento $\theta$ coerente col quadrante.
\item Per $z^n$, usa $r^n e^{in\theta}$.
\item Per $w^n=z$, usa $w_k=r^{1/n}e^{i(\theta+2k\pi)/n}$, $k=0,\ldots,n-1$.
\item Elenca tutte le radici e controlla che siano equidistanti sulla circonferenza di raggio $r^{1/n}$.
\end{enumerate}
\end{proceduraesame}
\begin{errorefrequente}{argomento}
L'argomento è multivalore modulo $2\pi$; l'argomento principale è una scelta di ramo. Confondere i due concetti fa perdere radici o produce logaritmi complessi errati.
\end{errorefrequente}

\section{Forma algebrica e operazioni}\label{forma-algebrica-e-operazioni}

Nelle formule seguenti \(a,b,c,d\in\mathbb R\); \(i\) è l'unità immaginaria.

\[
\displaystyle
z=a+ib,
\qquad
i^2=-1,
\qquad
\operatorname{Re}z=a,
\qquad
\operatorname{Im}z=b.
\]

Per \(z=a+ib\) e \(w=c+id\):

\[
\begin{aligned}
z+w&=(a+c)+i(b+d),\\
z-w&=(a-c)+i(b-d),\\
zw&=(ac-bd)+i(ad+bc).
\end{aligned}
\]

\[
\displaystyle
\overline z=a-ib,
\qquad
|z|=\sqrt{a^2+b^2},
\qquad
z\overline z=|z|^2.
\]

\[
\begin{aligned}
\overline{z+w}&=\overline z+\overline w,\\
\overline{zw}&=\overline z\,\overline w,\\
|zw|&=|z|\,|w|,\\
|z+w|&\le |z|+|w|.
\end{aligned}
\]

Per \(z\ne0\):

\[
\displaystyle
\frac1z=\frac{\overline z}{|z|^2}
=\frac{a-ib}{a^2+b^2}.
\]

Per \(w\ne0\):

\[
\displaystyle
\frac zw
=
\frac{z\overline w}{|w|^2}
=
\frac{(ac+bd)+i(bc-ad)}{c^2+d^2}.
\]

Potenze di \(i\):

\begin{longtable}[]{@{}rc@{}}
\toprule\noalign{}
\(n\pmod 4\) & \(i^n\) \\
\midrule\noalign{}
\endhead
\bottomrule\noalign{}
\endlastfoot
\(0\) & \(1\) \\
\(1\) & \(i\) \\
\(2\) & \(-1\) \\
\(3\) & \(-i\) \\
\end{longtable}

Approfondimenti: \href{https://ingegnerismo.it/matematica/numeri-complessi/}{numeri complessi} e \href{https://ingegnerismo.it/matematica/modulo-numero-complesso/}{modulo di un numero complesso}.

\section{Forma trigonometrica ed esponenziale}\label{forma-trigonometrica-ed-esponenziale}

Per \(z=a+ib\ne0\):

\[
\displaystyle
\rho=|z|=\sqrt{a^2+b^2},
\qquad
z=\rho(\cos\theta+i\sin\theta)=\rho e^{i\theta}.
\]

\[
\displaystyle
\cos\theta=\frac a\rho,
\qquad
\sin\theta=\frac b\rho,
\qquad
\tan\theta=\frac ba\quad(a\ne0),
\]

\[
a=\rho\cos\theta,
\qquad
b=\rho\sin\theta.
\]

Per scegliere il quadrante corretto non basta in generale \(\arctan(b/a)\): si usa l'argomento determinato congiuntamente dai segni di \(a\) e \(b\), spesso indicato operativamente con \(\operatorname{atan2}(b,a)\).

\[
\displaystyle
e^{i\theta}=\cos\theta+i\sin\theta.
\]

\[
\displaystyle
\arg z=\{\theta+2k\pi:k\in\mathbb Z\}.
\]

Se \(z=\rho e^{i\theta}\) e \(w=\sigma e^{i\varphi}\):

\[
\begin{aligned}
zw&=\rho\sigma e^{i(\theta+\varphi)},\\
\frac zw&=\frac\rho\sigma e^{i(\theta-\varphi)},\qquad w\ne0,\\
\overline z&=\rho e^{-i\theta}.
\end{aligned}
\]

Approfondimenti: \href{https://ingegnerismo.it/matematica/argomento-numero-complesso/}{argomento di un numero complesso} e \href{https://ingegnerismo.it/matematica/formula-di-eulero/}{formula di Eulero}.

\section{Potenze e radici}\label{potenze-e-radici}

Formula di De Moivre:

\[
\displaystyle
(\cos\theta+i\sin\theta)^n
=
\cos(n\theta)+i\sin(n\theta),
\qquad n\in\mathbb Z,
\]

con base non nulla se \(n<0\).

Per \(z=\rho e^{i\theta}\ne0\):

\[
\displaystyle
z^n=\rho^n e^{in\theta},
\qquad n\in\mathbb Z.
\]

\[
\displaystyle
0^n=0,
\qquad n\ge1.
\]

Per \(n\ge1\), \(z=\rho e^{i\theta}\ne0\), l'equazione \(w^n=z\) possiede esattamente \(n\) radici distinte:

\[
\displaystyle
w_k=\sqrt[n]{\rho}\,
e^{i(\theta+2k\pi)/n},
\qquad
k=0,1,\ldots,n-1.
\]

Cambiare \(\theta\) con \(\theta+2m\pi\) non produce nuove radici oltre alle \(n\) già elencate: ne modifica soltanto l'ordine.

Caso nullo:

\[
\displaystyle
w^n=0
\quad\Longleftrightarrow\quad
w=0,
\qquad n\ge1.
\]

Convenzione sulle potenze nulle:

\[
\displaystyle
z^0=1\quad(z\ne0),
\qquad
0^0\text{ non definito}.
\]

Approfondimenti: \href{https://ingegnerismo.it/matematica/formula-di-de-moivre/}{formula di De Moivre} e \href{https://ingegnerismo.it/matematica/radici-complesse/}{radici complesse}.

\section{Teorema fondamentale dell'algebra}\label{teorema-fondamentale-dellalgebra}

Per ogni polinomio complesso non costante

\[
P(z)=a_nz^n+a_{n-1}z^{n-1}+\cdots+a_1z+a_0,
\qquad a_n\ne0,
\qquad n\ge1,
\]

esistono \(z_1,\ldots,z_n\in\mathbb C\), ripetuti secondo la molteplicità, tali che

\[
P(z)=a_n\prod_{k=1}^{n}(z-z_k).
\]

Quindi un polinomio di grado \(n\) ha esattamente \(n\) radici complesse contando le molteplicità. Se \(P\in\mathbb R[z]\):

\[
P(z_0)=0
\Longrightarrow
P(\overline{z_0})=0,
\]

perciò le radici non reali compaiono a coppie coniugate.

\section{Esponenziale e logaritmo complesso (estensione)}\label{esponenziale-e-logaritmo-complesso-estensione}

\[
\displaystyle
e^{a+ib}=e^a(\cos b+i\sin b).
\]

Per \(z\ne0\):

\[
\displaystyle
\operatorname{Arg}z\in(-\pi,\pi],
\qquad
\arg z
=
\{\operatorname{Arg}z+2k\pi:k\in\mathbb Z\}.
\]

Logaritmo multivalore. A differenza del logaritmo reale, \(\log z\) è un insieme infinito di valori:

\[
\displaystyle
\log z
=
\left\{
\ln|z|+i\bigl(\operatorname{Arg}z+2k\pi\bigr):
k\in\mathbb Z
\right\}.
\]

\[
\displaystyle
e^w=z
\quad\Longleftrightarrow\quad
w\in\log z,
\qquad z\ne0.
\]

Ramo principale analitico con taglio sull'asse reale non positivo:

\[
\displaystyle
z\in\mathbb C\setminus(-\infty,0],
\qquad
\operatorname{Arg}z\in(-\pi,\pi),
\qquad
\operatorname{Log}z
=
\ln|z|+i\operatorname{Arg}z.
\]

Approfondimento: \href{https://ingegnerismo.it/matematica/logaritmo-complesso/}{logaritmo complesso}.
